\section{Complete proofs for the solved problems}\label{app:proofs}

The proofs below are transcriptions of the reports in the repository (\code{reports/*.md}); they are
self-contained.

\subsection{Problem 21.32}\label{app:2132}

Throughout, $H$ is a finite group with $H'=G$; we identify $G$ with the subgroup $H'\le H$ and put $K:=C_H(G)$,
$Z:=Z(G)$.

\begin{lemma}\label{lem:2132basic}
(1) $H/G$ is abelian. (2) $K\trianglelefteq H$ and $H/K$ embeds in $\Aut(G)$ via conjugation; its image $Q:=H/K$
contains $\Inn(G)=GK/K$. (3) $K\cap G=Z$. (4) $[H,K]\le Z$; in particular $K'\le Z$, $K$ has class $\le 2$, $K/Z$ is
central in $H/Z$, and $K$ centralises $Z$.
\end{lemma}
\begin{proof}
(1) $G=H'$. (2) The kernel of the conjugation action $H\to\Aut(G)$ is $C_H(G)=K$. (3) $K\cap G=C_G(G)=Z$.
(4) $[H,K]\le H'\cap K=G\cap K=Z$ because $K\trianglelefteq H$; the rest follows.
\end{proof}

\begin{theorem}[necessary condition]\label{thm:2132B}
If $G\cong H'$ for a finite group $H$, then there is a subgroup $Q$ with $\Inn(G)\le Q\le\Aut(G)$ and
$Q'=\Inn(G)$. If $Z(G)=1$ the converse holds.
\end{theorem}
\begin{proof}
Put $\tilde H=H/Z$, $\tilde K=K/Z$, $\tilde G=G/Z$. By Lemma~\ref{lem:2132basic}(4) $\tilde K$ is central in
$\tilde H$, and $\tilde H/\tilde K\cong H/K=Q$. Since $H'=G$ we get $\tilde H'=\tilde G$, and
$\tilde G\cap\tilde K=(G\cap K)/Z=1$. The projection $\pi:\tilde H\to Q$ maps $\tilde H'$ onto $Q'$ with kernel
$\tilde H'\cap\tilde K=1$, so $\pi$ restricts to an isomorphism $\tilde G=\tilde H'\to Q'$; and $\pi(\tilde G)$
is the image of $G$ in $H/K$, i.e.\ $\Inn(G)$. Hence $Q'=\Inn(G)$.
If $Z(G)=1$ then $K\cap G=1$ and $[H,K]=1$, so $K\le Z(H)$ and $H'\cap K=1$; hence $H'\cong H'K/K=(H/K)'=Q'=\Inn(G)
\cong G$, and conversely, if $Q\le\Aut(G)$ satisfies $\Inn(G)\le Q$ and $Q'=\Inn(G)\cong G$, then $H:=Q$ works.
\end{proof}

\begin{theorem}[witness bound]\label{thm:2132A}
Let $e=2\exp(Z)$ and $m=\lceil\log_2|G|\rceil$. If $G\cong H'$ for some finite $H$, then $G\cong (H^*)'$ for a finite
group $H^*$ with $|H^*|\le|\Aut(G)|\cdot|Z|^2\cdot e^{2m}$.
\end{theorem}

\begin{proof}
\emph{Step 1 (bounded rank).} $G=H'$ is generated by commutators $[a,b]$ ($a,b\in H$), and every generating set of a
group of order $|G|$ contains a generating subset of size $\le\log_2|G|$ (a chain of proper subgroups has length
$\le\log_2|G|$). Choose $a_1,b_1,\dots,a_m,b_m\in H$ with $G=\gen{[a_i,b_i]}$ and let
$H_0=\gen{G,a_1,b_1,\dots,a_m,b_m}$. Then $G\le H_0'\le H'=G$, so $H_0'=G$, and $H_0/G$ is abelian, generated by
$2m$ elements. Replacing $H$ by $H_0$ we assume from now on that $H/G$ is abelian and generated by $2m$ elements.
Then $K/Z\cong KG/G\le H/G$ is abelian and generated by $\le 2m$ elements (subgroups of a $d$-generated finite
abelian group are $d$-generated).

\emph{Step 2 (a large central subgroup).} For $k,k'\in K$ we have $[k,k']\in K'\le Z$ and $K$ has class $\le 2$, so
$(kk')^e=k^ek'^e[k',k]^{e(e-1)/2}$; since $e$ is even, $e(e-1)/2=(e/2)(e-1)$ is a multiple of $\exp(Z)$ and the
last factor is trivial. Hence $k\mapsto k^e$ is an endomorphism of $K$ and $C:=K^e=\{k^e:k\in K\}$ is a
characteristic subgroup of $K$, so $C\trianglelefteq H$. For $h\in H$ and $k\in K$, $[h,k]\in Z$ is centralised by
$k$, so $[h,k^e]=[h,k]^e=1$ (using $[h,kk']=[h,k'][h,k]^{k'}$ repeatedly). Thus $C\le Z(H)$. Furthermore $K/C$ has
exponent dividing $e$, and $K/CZ$ is an abelian group generated by $\le 2m$ elements of exponent dividing $e$, so
$|K:CZ|\le e^{2m}$ and
\begin{equation}\label{eq:2132bound}
|H/C|=|H:K|\cdot|K:CZ|\cdot|CZ:C|\le|\Aut(G)|\cdot e^{2m}\cdot|Z|.
\end{equation}

\emph{Step 3 (replacing $C$ by $C\cap G$).} Put $\bar H:=H/C$ and $C_0:=C\cap G=C\cap Z$. $H$ is a central extension
$1\to C\to H\to\bar H\to 1$ with class $\xi\in H^2(\bar H,C)$ (trivial action). We use two standard facts about a
central extension $1\to M\to E\to Q\to 1$ with class $\eta$:
\begin{itemize}
\item[(F1)] (Hopf) the five-term exact sequence $H_2(E)\to H_2(Q)\to M\to E^{ab}\to Q^{ab}\to 0$ shows that the
image of the Hopf map $h_\eta:H_2(Q)\to M$ equals $M\cap E'$;
\item[(F2)] (universal coefficients) $0\to\mathrm{Ext}(Q^{ab},M)\to H^2(Q,M)\to\mathrm{Hom}(H_2(Q),M)\to 0$ is
exact, the second map sending $\eta$ to $h_\eta$; the sequence is natural in $M$.
\end{itemize}
By (F1), $h_\xi$ has image $C\cap H'=C_0$. Let $\iota:C_0\hookrightarrow C$. By (F2) with coefficient group
$C_0$ there is $\xi^*\in H^2(\bar H,C_0)$ with $h_{\xi^*}=h_\xi$ (as a map into $C_0$). Let $H^*$ be a central
extension of $\bar H$ by $C_0$ with class $\xi^*$; then $|H^*|=|\bar H|\cdot|C_0|$. By naturality
$h_{\iota_*\xi^*}=\iota\circ h_{\xi^*}=h_\xi$, so $\xi-\iota_*\xi^*$ has trivial Hopf map, i.e.\ lies in
$\mathrm{Ext}(\bar H^{ab},C)$; let $A$ be a central extension of $\bar H$ by $C$ with this class. By (F1),
$A'\cap C=1$, so $A'\to\bar H'$ is an isomorphism.

The extension with class $\iota_*\xi^*$ is the pushout $E_1:=(H^*\times C)/\{(c,c^{-1}):c\in C_0\}$; the map
$H^*\to E_1$ is injective, $C$ is central in $E_1$, $E_1=H^*C$, and hence $E_1'=(H^*)'$.

Now $\xi=\iota_*\xi^*+(\xi-\iota_*\xi^*)$ is a Baer sum: $H\cong F/\Delta$ where $F:=E_1\times_{\bar H}A$ is the fibre
product over $\bar H$ (a central extension of $\bar H$ by $C\times C$) and $\Delta=\{(c,c^{-1}):c\in C\}$.

\emph{Claim:} $F'=E_1'\times_{\bar H'}A'\cong E_1'$. Indeed $F'\subseteq E_1'\times_{\bar H'}A':=\{(x,y)\in E_1'\times
A':x,y\text{ have the same image in }\bar H'\}$. By naturality of the Hopf map under the two projections, the Hopf
map of $F$ is $(h_{\iota_*\xi^*},h_{\xi-\iota_*\xi^*})=(h_\xi,0)$, with image $C_0\times 1$; by (F1),
$F'\cap(C\times C)=C_0\times 1$. As $F'$ maps onto $\bar H'$, $|F'|=|\bar H'|\cdot|C_0|$. On the other hand
$|E_1'\times_{\bar H'}A'|=\sum_{z\in\bar H'}|E_1'\cap\pi^{-1}(z)|\cdot|A'\cap\pi^{-1}(z)|=|\bar H'|\cdot|E_1'\cap C|
\cdot|A'\cap C|=|\bar H'|\cdot|C_0|\cdot 1$. Hence $F'=E_1'\times_{\bar H'}A'$, and the projection to the first
coordinate is an isomorphism onto $E_1'$ (injective because $A'\cap C=1$, surjective because $A'\to\bar H'$ is
onto).

Finally $F'\cap\Delta=\{(c,c^{-1}):c\in C_0,\ c^{-1}=1\}=1$, so $H'=(F/\Delta)'=F'\Delta/\Delta\cong F'\cong E_1'=
(H^*)'$. Therefore $(H^*)'\cong H'=G$ and, by \eqref{eq:2132bound},
$|H^*|=|\bar H|\cdot|C_0|\le|\Aut(G)|\cdot e^{2m}\cdot|Z|\cdot|Z|$.
\end{proof}

\begin{remark}
The witness has order $|G|^{O(\log|G|)}$ (as $|\Aut G|\le|G|^{\log_2|G|}$), and ``$(H^*)'\cong G$'' is verifiable
in time quasi-polynomial in $|G|$ (group isomorphism for groups of order $N$ is decidable in time
$N^{O(\log N)}$); so the decision problem is in nondeterministic quasi-polynomial time. For centreless $G$
Theorem~\ref{thm:2132B} gives a much better algorithm inside $\Aut(G)$. Examples: $\Aut(S_3)=S_3$ has derived
subgroup $C_3\ne\Inn(S_3)$, and no subgroup of $\Aut(D_8)\cong D_8$ containing $\Inn(D_8)\cong C_2^2$ has derived
subgroup $C_2^2$; so $S_3$ and $D_8$ are not derived subgroups of finite groups (confirmed by brute force in
\code{work/p21\_32/}).
\end{remark}

\subsection{Problem 16.60}\label{app:1660}

\begin{theorem}
For every finite group $G$ and every $\alpha\in\Aut(G)$, $|S_\alpha|\le T(G)$.
\end{theorem}
\begin{proof}
\emph{Step 1.} Put $C:=C_G(\alpha^2)=\{g:\alpha^2(g)=g\}$. If $g\in S_\alpha$ then $\alpha^2(g)=\alpha(g^{-1})=
\alpha(g)^{-1}=g$, so $S_\alpha\subseteq C$. Since $\alpha$ commutes with $\alpha^2$, $\alpha(C)=C$, and
$\sigma:=\alpha|_C\in\Aut(C)$ satisfies $\sigma^2=\mathrm{id}$. Clearly $S_\alpha=\{g\in C:\sigma(g)=g^{-1}\}=:
S_\sigma(C)$. It suffices to prove (A) $|S_\sigma(C)|\le T(C)$ for every finite group $C$ and involutory
$\sigma\in\Aut(C)$, and (B) $T(C)\le T(G)$ for subgroups $C\le G$.

\emph{Step 2 (B).} For $\theta\in\Irr(C)$ the induced character $\theta^G$ is a non-zero character, so
$m(\theta):=\sum_{\chi\in\Irr(G)}\langle\chi,\theta^G\rangle\ge 1$. By Frobenius reciprocity,
$T(G)=\sum_\chi\chi(1)=\sum_\chi\sum_\theta\langle\chi_C,\theta\rangle\theta(1)=\sum_\theta\theta(1)m(\theta)\ge
\sum_\theta\theta(1)=T(C)$.

\emph{Step 3 (A).} For $\theta\in\Irr(C)$ put $\varepsilon_\sigma(\theta):=\frac1{|C|}\sum_{g\in C}\theta(g\sigma(g))$.
(i) By the second orthogonality relation $\sum_\theta\theta(1)\theta(x)=|C|[x=1]$,
$\sum_\theta\varepsilon_\sigma(\theta)\theta(1)=\#\{g\in C:g\sigma(g)=1\}=|S_\sigma(C)|$.
(ii) $\varepsilon_\sigma(\theta)\in\{1,0,-1\}$: let $\rho:C\to\GL(V)$ afford $\theta$ and $\rho^\sigma:=\rho\circ\sigma$.
For linear maps $A,B$ on $V$, $\operatorname{tr}(AB)=\operatorname{tr}(\tau\circ(A\otimes B))$, where $\tau$ is the
flip on $V\otimes V$. Hence $\varepsilon_\sigma(\theta)=\operatorname{tr}(\tau\circ P)$ where
$P=\frac1{|C|}\sum_g(\rho\otimes\rho^\sigma)(g)$ is the projection onto the invariants $W$ of $\rho\otimes\rho^\sigma$.
Since $\sigma^2=1$, $\tau\circ(\rho\otimes\rho^\sigma)(g)=(\rho\otimes\rho^\sigma)(\sigma(g))\circ\tau$, so $\tau$
maps $W$ into itself and $\varepsilon_\sigma(\theta)=\operatorname{tr}(\tau|_W)$, an integer of absolute value
$\le\dim W=\dim\mathrm{Hom}_C(V^*,V^\sigma)=\langle\bar\theta,\theta^\sigma\rangle\le 1$.
(iii) From (i) and (ii), $|S_\sigma(C)|=\sum_\theta\varepsilon_\sigma(\theta)\theta(1)\le\sum_\theta\theta(1)=T(C)$.
\end{proof}
Step 3 is the twisted Frobenius--Schur theorem of Kawanaka and Matsuyama~\cite{KM}; Step 1 is the only new
ingredient. A numerical check for all groups of order $\le 200$ and all automorphisms is in \code{work/p16\_60/}.

\subsection{Problem 2.78}\label{app:278}

For a finite group $G$ let $\Ord(G)$ be the set of orders of subgroups of $G$, $f(G)=\#\{d:G\text{ has a non-normal
subgroup of order }d\}$ and $n(G)=|\Ord(G)|-f(G)$ (the number of orders $d$ all of whose subgroups are normal).

\begin{lemma}[sections]\label{lem:278A}
If $S=H/N$ with $N\trianglelefteq H\le G$, then $f(G)\ge f(S)$. For non-abelian simple $S$, $f(S)=|\Ord(S)|-2$.
\end{lemma}
\begin{proof}
If $K/N$ ($N\le K\le H$) is non-normal in $S$ then $K$ is not normal in $G$ (else $K\trianglelefteq H$ and
$K/N\trianglelefteq S$), and distinct $|K/N|$ give distinct $|K|=|N||K/N|$. For simple $S$ every proper non-trivial
subgroup is non-normal.
\end{proof}

\begin{lemma}[coprime direct products]\label{lem:278B}
If $\gcd(|G|,|H|)=1$ then $f(G\times H)=|\Ord(G)||\Ord(H)|-n(G)n(H)$.
\end{lemma}
\begin{proof}
Every subgroup of $G\times H$ is $A\times B$ with $A\le G$, $B\le H$, normal iff $A\trianglelefteq G$ and
$B\trianglelefteq H$; distinct pairs $(|A|,|B|)$ give distinct products. So the order $ab$ carries a non-normal
subgroup iff some subgroup of order $a$ is non-normal in $G$ or some subgroup of order $b$ is non-normal in $H$.
\end{proof}
In particular, for $G=A_5$ ($\Ord(A_5)=\{1,2,3,4,5,6,10,12,60\}$, $n(A_5)=2$): $f(A_5\times H)=9|\Ord(H)|-2n(H)$
whenever $\gcd(|H|,30)=1$.

\begin{lemma}[a family of $p$-groups]\label{lem:278C}
Let $p$ be an odd prime, $n\ge 3$, $r\ge 0$, $M=M_{p^n}=\gen{x,y\mid x^{p^{n-1}}=y^p=1,\ x^y=x^{1+p^{n-2}}}$ and
$P=M\times C_{p^r}$. Then $|\Ord(P)|=n+r+1$ and $f(P)=r+1$: the non-normal subgroups of $P$ have exactly the orders
$p,p^2,\dots,p^{r+1}$.
\end{lemma}
\begin{proof}
$P$ is a $p$-group of order $p^{n+r}$, so all orders $p^i$ ($0\le i\le n+r$) occur. Put $z=x^{p^{n-2}}$; then
$M'=\gen z$ has order $p$, $Z(M)=\gen{x^p}$, and $C=C_{p^r}$ is central in $P$. For $U\le P$ with projection
$V=\mathrm{pr}_M(U)$ we have $[U,P]=[V,M]\times 1$, which is $1$ if $V\le Z(M)$ and $\gen z\times 1$ otherwise; so
$U\trianglelefteq P$ iff $V\le Z(M)$ or $(z,1)\in U$.
(i) The subgroups $U=\gen{(y,1)}\times\gen{(1,c)}$ with $c\in C$ of order $p^j$ ($0\le j\le r$) satisfy
$V=\gen y\not\le Z(M)$ and contain no element $(z,1)$ (their elements are $(y^i,c^k)$ and $z\notin\gen y$), so
they are non-normal of order $p^{j+1}$.
(ii) Let $U$ be non-normal, so $V\not\le Z(M)$ and $(z,1)\notin U$. $V$ is abelian (otherwise $U$ would contain a
commutator $([m,m'],1)=(z^j,1)\ne 1$), so $V\le C_M(v)$ for some $v\in V\setminus Z(M)$, and $C_M(v)=\gen{v,Z(M)}$
has order $p^{n-1}$ (the centraliser of a non-central element has index $p$ since $M/Z(M)\cong C_p^2$). If
$v=x^iy^j$ with $p\nmid i$ then $v^p=x^{ip}$ ($p$ odd, class 2, $|M'|=p$), so $v$ has order $p^{n-1}$ and
$C_M(v)=\gen v$ is cyclic with unique subgroup of order $p$ equal to $\gen z$; if $p\mid i$ then $v=x^{pi'}y^j$ with $p\nmid j$, and replacing $v$ by a power generating the same subgroup we may
assume $v=yx^{pi'}$; then $C_M(v)=\gen y\times\gen{x^p}$ with $z\in\gen{x^p}$. In both cases a subgroup $V_0\le C_M(v)$ with $z\notin V_0$
meets the cyclic factor $\gen{x^p}$ (resp.\ $\gen v$) trivially, hence $|V_0|\le p$. Now $U\cap(M\times 1)=V_0\times 1$
with $z\notin V_0$, so $|V_0|\le p$ and $|U|=|V_0|\cdot|\mathrm{pr}_C(U)|\le p^{r+1}$.
\end{proof}
(The lemma was also checked by computer for $p=7,11$, $n\le 4$, $r\le 2$.)

\begin{theorem}\label{thm:278main}
Every $k\ge 86$ is $f(G)$ for a non-soluble non-simple group $G$. Hence the soluble numbers and the absolutely
simple numbers form finite sets (contained in $\{0,\dots,85\}$).
\end{theorem}
\begin{proof}
The numbers $7n+9$ for $n=3,\dots,11$ are $30,37,44,51,58,65,72,79,86$, which represent all residue classes modulo 9
and are $\le 86$. Hence every $k\ge 86$ can be written as $k=7n+9r+9$ with $3\le n\le 11$ and $r\ge 0$. Take a prime
$p\ge 7$ and $P=M_{p^n}\times C_{p^r}$. By Lemmas~\ref{lem:278B} and \ref{lem:278C},
$f(A_5\times P)=9(n+r+1)-2n=7n+9r+9=k$, and $A_5\times P$ is non-soluble and not simple.
\end{proof}

\begin{theorem}\label{thm:278small}
(a) $f(G)\ge 7$ for every non-soluble $G$; (b) $f(G)=7$ iff $G\cong A_5$; (c) no non-soluble group has $f(G)=8$;
(d) $f(S)\ge 9$ for every non-abelian simple $S\ne A_5$.
\end{theorem}
\begin{proof}
(d) Every non-abelian simple $S$ has a minimal non-soluble subgroup $U$, and $U/\Phi(U)$ is a minimal simple group $T$
(Thompson~\cite{Thompson}); by Lemma~\ref{lem:278A}, $f(S)\ge f(T)$. We check $f(T)\ge 9$, i.e.\ $|\Ord(T)|\ge 11$,
for every minimal simple $T\ne A_5$: for $T=\PSL(2,2^p)$ ($p\ge 3$ prime) $\Ord(T)$ contains the $p+1$ powers of 2 up to
$2^p$, the divisors of $2^p\pm1$ and their doubles (cyclic and dihedral subgroups), the Borel order $2^p(2^p-1)$ and
$|T|$; for $p=3$ this gives $|\Ord|=12$, and for $p\ge 5$ at least $6+6$. For $T=\PSL(2,3^p)$ ($p$ odd prime) similarly
$|\Ord(\PSL(2,27))|=14$ and more for larger $p$. For $T=\PSL(2,p)$, $p\ge 7$ prime, $\Ord(T)$ contains
$1,2,3,4,6,12,p,(p-1)/2,(p+1)/2,p-1,p+1,p(p-1)/2,|T|$ (Klein four groups, $S_3$, $A_4$, cyclic and dihedral
subgroups, the Borel subgroup); these 13 numbers are distinct unless $p\in\{7,11,13,23\}$; in those cases
$\Ord(\PSL(2,7))=\{1,2,3,4,6,7,8,12,21,24,168\}$ (using $S_4\le\PSL(2,7)$),
$\Ord(\PSL(2,11))\supseteq\{1,2,3,4,5,6,10,11,12,55,60,660\}$ ($A_5\le\PSL(2,11)$),
$\Ord(\PSL(2,13))\supseteq\{1,2,3,4,6,7,12,13,14,26,39,78,1092\}$ and
$\Ord(\PSL(2,23))\supseteq\{1,2,3,4,6,8,11,12,22,23,24,253,6072\}$, so $|\Ord|\ge 11$ throughout. For $T=\Sz(2^p)$ there are $2p+1$
two-power orders, three odd torus orders and $|T|$, so $|\Ord|\ge 2p+5\ge 11$. For $T=\PSL(3,3)$, $|\Ord|=23$.
So $f(S)\ge 9$ unless the minimal simple section is $A_5$. If $U/\Phi(U)\cong A_5$ with $\Phi(U)\ne 1$, then
besides the 7 orders $|\Phi(U)|d$ ($d\in D:=\{2,3,4,5,6,10,12\}$) the subgroups $\Phi(U)$ and $U$ give two further
non-normal orders, so $f(S)\ge 9$. If $A_5\cong U\le S\ne A_5$, then $\Ord(S)\supseteq\Ord(A_5)\cup\{60\}$ gives
$f(S)\ge 8$, and a Sylow 2-subgroup of order $\ge 8$ gives a ninth order unless $|S|_2=4$; then $S\cong\PSL(2,q)$ with
$q\equiv\pm3\pmod 8$ (Gorenstein--Walter), $q\ge 11$, and $q\in\Ord(S)$ is a new order.

(a) follows from Lemma~\ref{lem:278A}, (d) and $f(A_5)=7$.

(b), (c). Let $G$ be non-soluble with $f(G)\in\{7,8\}$. By (d) and Lemma~\ref{lem:278A} every non-abelian
composition factor of $G$ is $A_5$; fix $H\le G$, $N\trianglelefteq H$ with $H/N\cong A_5$. The 7 orders $|N|d$
($d\in D$) are orders of non-normal subgroups of $G$; at most one further order $e$ is available.

\emph{$N=1$.} For $p\in\{2,3,5\}$ a Sylow $p$-subgroup $Q$ of $H$ is non-normal in $G$ (its image in $H/N$ is a
non-normal Sylow subgroup of $A_5$), and $|Q|=|A_5|_p|N|_p$. If $|Q|=|N|d$ with $d\in D$ then $|N|_{p'}d=|A_5|_p$,
forcing $d=|A_5|_p$ and $N$ a $p$-group. If $N\ne 1$, $N$ is a $p$-group for at most one of the three primes, so
at least two of the numbers $4|N|_2,3|N|_3,5|N|_5$ would have to equal the single extra order $e$; but no two of
them are equal. Hence $N=1$ and $A_5\cong H\le G$.

\emph{Minimal normal subgroups.} Let $M$ be a minimal normal subgroup of $G$. If $M\cap H=1$ then for each
non-normal $K\le H$ the subgroup $KM$ is non-normal in $G$ ($KM\trianglelefteq G$ would give
$KM\cap H=K(M\cap H)=K\trianglelefteq H$), of order $|K||M|$ with $|M|\ge 2$; the two orders $10|M|\ge 20$ and
$12|M|\ge 24$ lie outside $D$ and are distinct, so two extra orders would be needed. Hence $M\ge H$, so
$M\cong A_5^k$ (all non-abelian composition factors are $A_5$); $k\ge 2$ is impossible since $A_5\times A_5$ has
non-normal subgroups of the nine orders $2,3,4,5,6,9,10,12,25$. Thus $M=A_5$ is the unique minimal normal
subgroup, $C_G(M)=1$ and $G$ embeds in $\Aut(A_5)=S_5$. As $f(S_5)=10$, $G=A_5$ and $f(G)=7$. So no non-soluble
group has $f=8$, and $f=7$ forces $G\cong A_5$.
\end{proof}

\begin{remark}
The values $f(G)$ attained by non-soluble groups of order $\le 1000$ are $7,9,10,12,\dots,18,20,\dots,23$; small simple
groups give $f(\PSL(2,7))=9$, $f(\PSL(2,8))=f(\PSL(2,11))=10$, $f(\PSL(2,13))=11$, \dots, $f(M_{11})=24$, and
$f(\PSL(2,q))$ for $q\le 169$ takes the values $7$, $9$--$23$, $25$, $26$, $28$, $32$, $35$, $42$; together with
$f(M_{11})=24$ every integer in $\{7,9,\dots,26\}\setminus\{8\}$ is a simple number. Whether composite
non-soluble numbers exist (i.e.\ whether some $k\ge 9$ is avoided by all simple groups) remains open.
\end{remark}

\subsection{Problem 18.46}\label{app:1846}

\begin{theorem}
Let $H$ be a finite group containing a subgroup $D\cong D_8$ such that every element of $D$ is a square in $H$. Then
$|H|\ge 128$; the wreath product $H=D_8\wr C_2=(D_8\times D_8)\rtimes\gen\sigma$ attains the bound, the diagonal
subgroup $\{(g,g)\}\cong D_8$ consisting of squares since $((g,1)\sigma)^2=(g,1)(1,g)=(g,g)$.
\end{theorem}
\begin{proof}
If $|H|<128$ and $D_8\le H$ then $8\mid|H|$, so $|H|\in\{8,16,\dots,120\}$; all groups of these orders are in the
SmallGroups library. Two independent GAP computations (\code{work/p18\_46/check.g} using
\code{IsomorphicSubgroups(H,D8)}, and \code{verify\_d8.g} using the full subgroup lattice and \code{IdGroup}) show
that no group of order $8,16,\dots,120$ has a subgroup $S\cong D_8$ with $S\subseteq\{h^2:h\in H\}$. Exactly 14 of the
2328 groups of order 128 (SmallGroup ids 134, 136, 138, 139, 140, 141, 144, 146, 928--933) do, among them
$D_8\wr C_2=\mathrm{SmallGroup}(128,928)$.
\end{proof}
The same computation gives the minimal $|H|$ for the sixteen groups of order $\le 12$ listed in the report (the remaining groups of order $\le 12$ have odd order, where $H=G$ works since every element is a square, or are
$C_{10}$, where $H=C_{20}$ works);
$D_8$ is the only one of them needing $|H|>|G|^2$ (e.g.\ $C_2^2$: 16, $S_3$: 36, $Q_8$: 48, $C_2^3$: 64, $A_4$: 24, $D_{12}$: 144). Hence
Klyachko's estimate $|H|\le 2|G|^2$ cannot be improved to $c|G|^2$ with $c<2$.

\subsection{Problem 13.19}\label{app:1319}

\paragraph{Reformulation.} If $Q\le G_1\times\dots\times G_n$ and $N_i=Q\cap\ker\pi_i$, then $\bigcap N_i=1$, and
``$H$ projects onto $G_i$'' means $HN_i=Q$. Conversely, a group $Q$ with normal subgroups $N_1,\dots,N_n$ with
$\bigcap N_i=1$ and a normal subgroup $H$ with $HN_i=Q$ for all $i$ gives such a configuration with $G_i=Q/N_i$.
One checks $\gamma_n(Q)\le H\cdot\bigcap_i N_i=H$ (a commutator of weight $n$ has, for each $i$, a factor that can be
taken in $N_i$ modulo $H$), so $Q/H$ has class $\le n-1$; in particular $n\ge 3$ is needed for a non-regular
2-group $Q/H$.

\paragraph{The explicit example ($n=3$, $p=2$).} Let $V=\F_2^2$ with the alternating form $\gamma(v,w)=v_1w_2+v_2w_1$
and the bilinear map $c_1(v,w)=v_1w_2$, and let $A=\F_2^3$ with basis $u_{12},u_{13},u_{23}$; put $w_{11}=u_{12}$,
$w_{22}=u_{23}$, $w_{33}=u_{13}$, $w_{ij}=w_{ji}=u_{ij}$ ($i<j$). On $Q=V^3\times A$ define
$(x,a)(y,b)=(x+y,a+b+c(x,y))$ with $c(x,y)=\sum_i c_1(x_i,y_i)w_{ii}+\sum_{i<j}\gamma(x_i,y_j)w_{ij}$. Since $c$ is
bilinear this is a group of order $2^9$ (a central extension of $V^3$ by $A$), generated by involutions $e_i,f_i$
(the standard basis of the $i$-th copy of $V$) and central involutions $u_{12},u_{13},u_{23}$, with
$[e_i,f_i]=w_{ii}$, $[e_i,f_j]=[f_i,e_j]=u_{ij}$ ($i\ne j$), all other commutators trivial. Put
\begin{gather*} N_i=\{(x,a):x_j=0\ (j\ne i),\ a\in\gen{u_{ij},u_{ik}}\}\quad(\{i,j,k\}=\{1,2,3\}),\\
\hat H=\{(x,a):x_1+x_2+x_3=0\},\qquad A_0=\gen{u_{12}+u_{13},u_{13}+u_{23}}. \end{gather*}
Then: (1) $N_i$ is a normal subgroup of order 16, $N_i\cap N_j=\gen{u_{ij}}$ and $N_1\cap N_2\cap N_3=1$;
(2) $\hat H$ is a subgroup of order $2^7$ containing $A$, and $\hat H/A_0$ is elementary abelian and central in
$Q/A_0$ (modulo $A_0$ all $w_{ij}$ become equal, so commutators and squares of elements with $\sum x_i=0$ vanish);
(3) choosing a complement $H/A_0$ of $A/A_0$ in $\hat H/A_0$ gives a normal subgroup $H$ of order $2^6$ with
$H\cap A=A_0$; (4) $HN_i=Q$ for all $i$ (as $H\cap N_i=A_0\cap N_i$ has order 2); (5) $Q/H$ is a non-abelian group of
order 8, namely $D_8$ (the commutator form induced by $\gamma(\sum x_i,\sum y_i)$ is non-degenerate and the squaring
map is isotropic). With $G_i=Q/N_i$ (of order 32) this is a configuration as in the problem with $Q/H\cong D_8$,
which is not regular. All of (1)--(5) were verified in GAP from the pc presentation
(\code{work/p13\_19/construct\_pc2.g}) with $H=\gen{e_1e_2u_{23},f_1f_2u_{23},e_2e_3u_{23},f_2f_3u_{23},u_{12}u_{23},
u_{13}u_{23}}$.

\paragraph{The general construction.}
\begin{lemma}[vanishing of the top difference]\label{lem:mobius}
Let $G$ be nilpotent of class $c$ and $w$ a word in $n>c$ blocks of variables. For a tuple $a$ of elements of $G$
(grouped in $n$ blocks) and $S\subseteq[n]$ let $a_S$ be $a$ with the entries of the blocks outside $S$ replaced by
$1$. If $w(a_S)=1$ for every proper $S\subsetneq[n]$, then $w(a)=1$.
\end{lemma}
\begin{proof}
Let $N$ be the free nilpotent group of class $c$ on the variables, $\sigma_S:N\to N$ the endomorphism killing the
variables outside $S$ (so $\sigma_S\sigma_T=\sigma_{S\cap T}$), and $M$ the normal closure of
$\{\sigma_S(w):S\subsetneq[n]\}$. Since $\sigma_T(M)\le M$, the $\sigma_T$ act on $\bar N=N/M$, and
$\sigma_S(\bar w)=1$ for all proper $S$. Suppose $\bar w\in\gamma_k(\bar N)$ with $k\le c$ and let
$A=\gamma_k(\bar N)/\gamma_{k+1}(\bar N)$, spanned by weight-$k$ commutators in the generators; for $T\subseteq[n]$
let $A_T$ be the span of those involving exactly the blocks in $T$. Then $\sigma_S$ acts on $A_T$ as the identity if
$T\subseteq S$ and as $0$ otherwise, and $\bar w\gamma_{k+1}=\sum_T a_T$ with $a_T\in A_T$, where only $|T|\le k<n$
occurs. Hence $0=\sigma_S(\bar w\gamma_{k+1})=\sum_{T\subseteq S}a_T$ for every proper $S$, and since
$\sum_{T\subseteq S\subsetneq[n]}(-1)^{n-|S|+1}=1$ for every proper $T$, we get
$0=\sum_{S\subsetneq[n]}(-1)^{n-|S|+1}\sigma_S(\bar w\gamma_{k+1})=\sum_T a_T=\bar w\gamma_{k+1}$. So
$\bar w\in\gamma_{k+1}(\bar N)$, and inductively $\bar w\in\gamma_{c+1}(\bar N)=1$. The homomorphism $N\to G$
sending the variables to $a$ maps $M$ to $1$ by hypothesis, hence factors through $\bar N$, and $w(a)$ is the image
of $\bar w=1$.
\end{proof}

\begin{theorem}\label{thm:1319gen}
Let $P$ be a finite $p$-group of class $c$ and $n=c+1$. Let $D=\prod_{\emptyset\ne S\subseteq[n]}P^{(S)}$; for a
generator $g$ of $P$ and $i\in[n]$ let $g^{[i]}\in D$ have $S$-component $g$ if $i\in S$ and $1$ otherwise; let $Q$ be
the subgroup generated by all $g^{[i]}$, $H=Q\cap\ker\pi_{[n]}$ and $N_i=Q\cap\bigcap_{S\not\ni i}\ker\pi_S$. Then
$Q/H\cong P$, $HN_i=Q$ for all $i$, and $N_1\cap\dots\cap N_n=1$; so $Q$, $H$ and $G_i=Q/N_i$ form a configuration
as in Problem 13.19 with $Q/H\cong P$.
\end{theorem}
\begin{proof}
$\pi_{[n]}(g^{[i]})=g$, so $\pi_{[n]}(Q)=P$ and $Q/H\cong P$. $N_i$ is normal (an intersection of $Q$ with kernels),
$g^{[i]}\in N_i$, so $\pi_{[n]}(N_i)=P$ and $HN_i=Q$. If $q\in\bigcap N_i=Q\cap P^{([n])}$ then $q$ is a word $w$ in
the tagged generators with $\pi_S(q)=w(a_S)=1$ for all proper $S$ ($a$ the tuple of generators grouped in $n$
blocks), so $\pi_{[n]}(q)=w(a)=1$ by Lemma~\ref{lem:mobius}, i.e.\ $q=1$.
\end{proof}
Taking $P=C_p\wr C_p$ (class $p$, irregular) gives a negative answer for every prime $p$. The construction was
checked in GAP for $P=D_8$ ($|Q|=2^{12}$) and $P=C_3\wr C_3$ ($|Q|=3^{32}$).
